\section{A constructive normal form for the full Cayley shuffle ideal}
\label{cay:sec:projection}

The earlier Cayley report determines the dimensions of two formal
presentations: the span of the unmultiplied Cayley rows, and the larger
shuffle ideal generated by them.  Its weight-seven search notes identify
an explicit reconstruction map for the latter quotient as a missing
computational step.  We provide such a map in every weight.  It gives
exact representatives without a relation matrix and produces a concrete
obstruction for the frozen $S_6$ formula after \emph{all} shuffle multiples
of Cayley identities have been included.

The shuffle-polynomial theorem and Eulerian idempotents used below are
classical \cite{cay:Radford,cay:Patras}.  The contribution to the inspected
project is the compatible endpoint projection, its combination with
those tools, the explicit reconstruction formula, and the resulting
weight-seven certificate.  Neither a period-independence assertion nor
the conjectured $S_6$ equality is established here.

\subsection{The algebra and its precise relation ideal}

Use the outer-letter-first convention and the five letters
\[
 z=\frac{dt}{t},\quad c=\frac{dt}{1-t},\quad
 a=\frac{i\,dt}{1-it},\quad b=-\frac{dt}{1+t},\quad
 \bar a=-\frac{i\,dt}{1+it}.
\]
Let $\widetilde{\mathcal A}$ be the rational shuffle algebra on these
letters, with empty word $1$ and shuffle product $\mathbin{\sqcup\!\sqcup}$.
Write $\mathcal A$ for its subalgebra spanned by $1$ and all words with
first letter different from $c$ and last letter different from $z$.
For such a word $w$, $H(w)$ is its convergent iterated integral from
zero to one.  Products written inside a word are concatenations unless
the shuffle symbol is displayed.

The Cayley involution is
\[
 C(v_1\cdots v_n)=\tau(v_n)\cdots\tau(v_1),\qquad
 \begin{array}{c|ccccc}
 v&z&c&a&b&\bar a\\ \hline
 \tau(v)&c-b&z+b&b-\bar a&b&b-a .
 \end{array}
\]
It preserves $\mathcal A$, is an automorphism for the shuffle product,
and commutes with the conjugation $K$ that exchanges $a,\bar a$ and
fixes $z,c,b$.  The change of variable $t\mapsto(1-t)/(1+t)$ proves
$H(Cw)=H(w)$ for admissible $w$.  Define
\begin{equation}\label{cay:eq:ideal}
 \mathcal J=\langle f-Cf:f\in\mathcal A\rangle_{\mathrm{shuffle}}.
\end{equation}
Thus $\mathcal J$ includes arbitrary products of convergent Cayley
relations with admissible words.  No stuffle, distribution, or other
functional equation is imposed in this definition.

\subsection{Endpoint projection that commutes with Cayley symmetry}

The ordinary projection that sets the divergent shuffle generators
$z,c$ to zero is unsuitable: it does not commute with $C$.  For example,
the ordinary projection of $Cz=c-b$ is $-b$, whereas $C$ applied to
the projection of $z$ is zero.  The needed correction is small but
essential.  Put
\begin{equation}\label{cay:eq:shifted-endpoints}
 x=z+\frac b2,\qquad y=c-\frac b2.
\end{equation}
Then $Cx=y$, $Cy=x$, and $Kx=x$, $Ky=y$.

\begin{lemma}[Equivariant endpoint projection]\label{cay:lem:regularization}
There is a unique shuffle-algebra retraction
$\mathcal R:\widetilde{\mathcal A}\to\mathcal A$ that is the identity
on $\mathcal A$ and sends $x,y$ to zero.  It commutes with $C,K$.
Equivalently, it substitutes $z=-b/2$, $c=b/2$ in
\emph{shuffle-polynomial coordinates}.

For a word $w$ of length $n$, let $l$ be its initial run of $c$'s and
$r$ its terminal run of $z$'s.  If $w[u:n-v]$ denotes deletion of
$u$ initial letters and $v$ terminal letters, then
\begin{equation}\label{cay:eq:regularization}
 \mathcal R(w)=
 \sum_{u=0}^{l}\sum_{v=0}^{r}
  \left(\frac b2-c\right)^u
  \mathbin{\sqcup\!\sqcup} w[u:n-v]
  \mathbin{\sqcup\!\sqcup}
  \left(-\frac b2-z\right)^v.
\end{equation}
The powers in this formula are concatenation powers of linear forms.
\end{lemma}

\begin{proof}
Order the alphabet as $z<a<b<\bar a<c$.  Every Lyndon word of length
at least two is admissible.  The shuffle-polynomial theorem
\cite{cay:Radford} and the
triangular Lyndon factorization therefore give
$\widetilde{\mathcal A}=\mathcal A[z,c]=\mathcal A[x,y]$ as polynomial
algebras.  Setting $x,y$ to zero defines the claimed retraction.  The
ideal $(x,y)$ is stable under both $C,K$, and both automorphisms preserve
$\mathcal A$; hence the retraction commutes with them.

For the formula, deletion of one initial $c$ and deletion of one terminal
$z$ are commuting derivations of the shuffle algebra.  Their Taylor
substitutions replace $c$ by $b/2$ and $z$ by $-b/2$.  The $u$th shuffle
power of a single linear form is $u!$ times its concatenation power,
which cancels the Taylor factorial.  This yields
\eqref{cay:eq:regularization}.  In particular every term left after
cancellation is admissible; this also follows from the range of the
retraction.
\end{proof}

This is a formal polynomial projection.  No value is assigned to a
divergent iterated integral, and no regularized analytic identity is
being assumed.  For example,
\[
 \mathcal R(az)=-za-\frac{ba+ab}{2},\qquad
 \mathcal R(c)=\frac b2,\qquad \mathcal R(z)=-\frac b2.
\]
The first formula also illustrates why a literal letter substitution
inside concatenation words would be incorrect.

\subsection{Explicit eigenvector lifts and the normal-form theorem}

For a nonempty word $w$, define the first Eulerian operator by
\begin{equation}\label{cay:eq:eulerian}
 e(w)=\sum_{k=1}^{|w|}\frac{(-1)^{k-1}}{k}
 \sum_{\substack{w=w_1\cdots w_k\\ |w_j|>0}}
  w_1\mathbin{\sqcup\!\sqcup}\cdots
  \mathbin{\sqcup\!\sqcup}w_k,\qquad e(1)=0.
\end{equation}
The inner sum is over consecutive nonempty blocks.  This is the
convolution logarithm of the identity for shuffle multiplication and
deconcatenation.  It annihilates a shuffle product of two positive-weight
elements and satisfies
\begin{equation}\label{cay:eq:eulerian-properties}
 e^2=e,\qquad e(f)-f\in
   (\widetilde{\mathcal A}_{>0})^2
 \quad(f\in\widetilde{\mathcal A}_{>0}).
\end{equation}
These are the standard Eulerian-idempotent properties in a connected
commutative Hopf algebra \cite{cay:Patras}.  The congruence is immediate
from \eqref{cay:eq:eulerian}; together with annihilation of products it
also proves the displayed idempotence.  Reversing the consecutive blocks
and using commutativity of shuffle shows directly that $e$ commutes
with $C$; it also commutes with $K$.

Set
\begin{equation}\label{cay:eq:lifts}
 E=\mathcal R e,\qquad
 E_+(w)=\frac{E(w)+C E(w)}2.
\end{equation}
All values of $E,E_+$ are admissible word polynomials.  The following
formula contains only finite sums of explicitly given rational words.

\begin{theorem}[Constructive full-ideal projector]\label{cay:thm:projector}
Define $\Pi(1)=1$ and, for every nonempty word,
\begin{equation}\label{cay:eq:projection}
 \boxed{\displaystyle
 \Pi(w)=\sum_{k=1}^{|w|}\frac1{k!}
   \sum_{\substack{w=w_1\cdots w_k\\ |w_j|>0}}
    E_+(w_1)\mathbin{\sqcup\!\sqcup}\cdots
     \mathbin{\sqcup\!\sqcup}E_+(w_k).}
\end{equation}
Then $\Pi:\widetilde{\mathcal A}\to\mathcal A$ is a graded
shuffle-algebra homomorphism.  Its restriction to $\mathcal A$ satisfies
\begin{equation}\label{cay:eq:projection-properties}
 \Pi^2=\Pi,\qquad \Pi C=C\Pi=\Pi,\qquad \Pi K=K\Pi,
 \qquad \ker(\Pi|_{\mathcal A})=\mathcal J.
\end{equation}
Consequently $\Pi(f)=\Pi(g)$ if and only if $f-g\in\mathcal J$ for
$f,g\in\mathcal A$.  In particular, every admissible word has the
proved period identity
\begin{equation}\label{cay:eq:period-normal-form}
 H(w)=H(\Pi(w)).
\end{equation}
\end{theorem}

\begin{proof}
First restrict $E$ to $\mathcal A$.  It annihilates products of two
positive-weight elements, and $E(f)-f\in(\mathcal A_{>0})^2$ because
$\mathcal R$ preserves products and positive degree.  Hence $E^2=E$
on $\mathcal A$, and its image $U=E(\mathcal A)$ maps isomorphically to
the indecomposables
$\mathcal A_{>0}/(\mathcal A_{>0})^2$.  Lifts of any homogeneous basis
of those indecomposables freely generate the polynomial algebra
$\mathcal A$: the change from its Lyndon generators is degree triangular
with invertible linear part.  Thus multiplication gives an isomorphism
$\operatorname{Sym}(U)\simeq\mathcal A$.

Since $\mathcal R,e$ commute with $C,K$, the space $U$ is stable under
both.  Write $U=U_+\oplus U_-$ for its $C$ eigenspaces.  The algebra
map that fixes $U_+$ and kills $U_-$ is an idempotent map $P$ on
$\mathcal A$.  Its kernel is the ideal generated by $U_-$.  This ideal
equals $\mathcal J$: if $u\in U_-$, then $u-Cu=2u$; conversely $C$
acts as the identity after the negative generators have been killed,
so every $f-Cf$ belongs to that ideal.  The map $P$ commutes with $K$.

It remains to verify that the fully explicit formula is this map.
For every nonempty $w$, convolution exponentiation gives
\[
 w=\sum_{k=1}^{|w|}\frac1{k!}
   \sum_{w=w_1\cdots w_k}
    e(w_1)\mathbin{\sqcup\!\sqcup}\cdots
     \mathbin{\sqcup\!\sqcup}e(w_k).
\]
Applying $\mathcal R$ gives the same formula for $\mathcal R(w)$
with factors $E(w_j)$.  Those factors belong to $U$, including when
the blocks themselves are not admissible.  To see this last point,
write $\widetilde{\mathcal A}=\mathcal A[x,y]$.  In degree at least
two, every term of $w-\mathcal R(w)$ that involves $x$ or $y$ is a
product of positive-degree factors and is annihilated by $e$; in degree
one, applying $\mathcal R e$ kills $x,y$ directly.  Thus
$E(w)=E(\mathcal R(w))\in U$ in every degree.

Now apply the algebra map $P$.  On $U$, it acts by $(1+C)/2$, so the
result is exactly \eqref{cay:eq:projection}.  Therefore
$\Pi=P\mathcal R$, proving multiplicativity and every assertion in
\eqref{cay:eq:projection-properties}.  Finally, $w-\Pi(w)\in\mathcal J$
for admissible $w$.  The map $H$ annihilates this ideal by the convergent
Cayley identity and the shuffle product law, proving
\eqref{cay:eq:period-normal-form}.
\end{proof}

For example, the three admissible weight-one letters give
\[
 \Pi(a)=\frac{a+b-\bar a}{2},\qquad
 \Pi(\bar a)=\frac{\bar a+b-a}{2},\qquad \Pi(b)=b.
\]
Evaluating the first identity recovers
$\Li_1(i)+\Li_1(-i)=\Li_1(-1)$, including its sign.  At higher weights
the projector automatically includes the products of the corresponding
Cayley-negative generators, which the raw average $(1+C)/2$ does not
annihilate.  The image is the same quotient whose weight-seven imaginary
dimension was previously proved to be $7518$.  Formula
\eqref{cay:eq:projection} supplies representatives in the original word
coordinates; it does not claim that the resulting sparse supports are
already a chosen $7518$-element basis.

\subsection{An exact obstruction for the frozen weight-seven target}

The series convention is
$\Li_{r,s}(u,v)=\sum_{n>m\ge1}u^nv^m/(n^rm^s)$.  Consequently
\[
 H(z^{r-1}a z^{s-1}a)=\Li_{r,s}(i,1),\qquad
 H(z^5a\bar a)=\Li_{6,1}(i,-1).
\]
Here $G=\beta(2)$ is Catalan's constant, and
\[
 S_6=\sum_{n=0}^{\infty}\frac{(-1)^nH_n}{(2n+1)^6},\qquad
 H_n=\sum_{j=1}^{n}\frac1j,\qquad H_0=0,
 \quad g_{r,s}=\operatorname{Im}\Li_{r,s}(i,1).
\]
Let $F$ be the following explicitly fixed formal lift of the conjectured
integer relation:
\begin{align}
 F={}&-179712000z^5aa+485683200z^5a\bar a
       -36864000z^3az^2a+401080320zaz^4a\notag\\
 &-\frac{47600087040}{61}z^6a
       +11750400(za\mathbin{\sqcup\!\sqcup}z^4c)\notag\\
 &+109347840(z^3a\mathbin{\sqcup\!\sqcup}z^2c)
       -971366400(z^5a\mathbin{\sqcup\!\sqcup}b).
 \label{cay:eq:target}
\end{align}
Indeed $\beta(7)=61\pi^7/184320$, $\Li_1(-1)=-\log2$, and the
harmonic parity split gives
\begin{align*}
 \operatorname{Im}H(F)={}&485683200S_6-665395200g_{6,1}
 -36864000g_{4,3}+401080320g_{2,5}\\
 &-258247\pi^7+11750400G\zeta(5)
 +109347840\beta(4)\zeta(3)+971366400\beta(6)\log2.
\end{align*}

For a word polynomial $f$, define its imaginary coordinate reduction
$\mathcal O(f)$ by replacing each conjugate pair $w,Kw$ by
$(f_w-f_{Kw})w$, choosing the lexicographically smaller representative
under the code order $z<c<a<b<\bar a$; conjugation-fixed words are
discarded.  This convention encodes imaginary parts without a factor
of $1/2$.

\begin{lemma}[An all-weight separating coordinate]\label{cay:lem:coordinate}
For an integer $w\ge3$, put
$\ell_w(f)=[z^{w-2}aa]\mathcal O(\Pi(f))$.
For positive integers $r,s$ with $r+s=w$,
\begin{equation}\label{cay:eq:coordinate-family}
 \ell_w(z^{r-1}az^{s-1}a)=\frac{\delta_{s,1}}4,
 \qquad
 \ell_w(z^{w-2}a\bar a)=-\frac14.
\end{equation}
Moreover $\ell_w$ vanishes on every word with at most one occurrence
of $a$ or $\bar a$.
\end{lemma}
\begin{proof}
Write $N=w-2>0$.  To contribute to $z^Naa$, a projected word on
$\{z,a,\bar a\}$ must retain every zero letter and both nonreal
letters.  A block consisting only of zero letters cannot do this:
its Eulerian image is zero if its length exceeds one, and
$\mathcal R(z)=-b/2$ otherwise.  Also $CE$ of any such block contains
no $z$, since $E$ introduces no $c$ and $\tau(z)=c-b$.
Thus only one-block and two-block terms of
\eqref{cay:eq:projection} can contribute, and each block containing
a zero letter must use the $E/2$ part of $E_+$.

For a block containing just one $a$, its component with all zeros
retained is the usual right-endpoint shuffle regularization
\[
 z^uaz^v\longmapsto
 (-1)^v\binom{u+v}{v}z^{u+v}a.
\]
Indeed setting the shuffle variable $z$ to zero kills every product
term in $e(z^uaz^v)-(z^uaz^v)$, because each such product has a
factor consisting only of zeros.  The displayed formula follows by
shuffle regularization, or by induction on $v$ using the shuffle
product with $z$.

For $z^{r-1}az^{s-1}a$, the sum of the two-block shuffle coefficients
at $z^Naa$ is therefore
\[
 \Sigma=2\sum_{j=0}^{s-1}(-1)^j
   \binom{r-1+j}{j}\binom{N}{r-1+j}
 =2\binom{N}{r-1}\sum_{j=0}^{s-1}(-1)^j\binom{s-1}{j}
 =2\delta_{s,1}.
\]
Convolution exponentiation before the Cayley projection says that
the one-block coefficient is $\delta_{s,1}-\Sigma/2$.
After projection, the one-block term is multiplied by $1/2$ and
the two-block sum by $1/8$.  Their total is
$\delta_{s,1}/2-\Sigma/8=\delta_{s,1}/4$.
There is no $z^N\bar a\bar a$ contribution: obtaining both
conjugate letters would require the $CE$ part of both blocks, and
at least one block contains a zero letter.

For $z^Na\bar a$, the only contributing cut is $z^Na\mid\bar a$.
The first factor contributes $z^Na/2$ and the second contributes
$-a/2$ from $C\bar a=b-a$.  The exponential factor $1/2!$ and
the coefficient $2$ of $z^Naa$ in $z^Na\mathbin{\sqcup\!\sqcup}a$
give $-1/4$.  Again the conjugate coordinate is absent.  Finally,
$e,\mathcal R,C$ never increase the number of nonreal letters.
Every term of \eqref{cay:eq:projection} consequently has at most
one such letter when the original word does, proving the last assertion.
\end{proof}

\begin{proposition}[Full Cayley-ideal separation]\label{cay:prop:separation}
The rational functional
\[
 \ell(f)=[z^5aa]\,\mathcal O(\Pi(f))
\]
annihilates the full ideal $\mathcal J$ and satisfies
\begin{equation}\label{cay:eq:separation}
 \ell(F)=-166348800\ne0.
\end{equation}
Thus the imaginary lift of the frozen target is not a consequence of
convergent Cayley symmetry and its complete shuffle-multiplicative
closure.
\end{proposition}

\begin{proof}
Annihilation follows from Theorem~\ref{cay:thm:projector}.
Lemma~\ref{cay:lem:coordinate} gives
\[
\begin{array}{c|ccccc}
 f&z^5aa&z^5a\bar a&z^3az^2a&zaz^4a&z^6a\\ \hline
 \ell(f)&1/4&-1/4&0&0&0 .
\end{array}
\]
Each word in the three shuffle products in \eqref{cay:eq:target}
has just one nonreal letter, so those products also have value zero
under $\ell$.  The table gives
$(-179712000-485683200)/4=-166348800$.
The supplied program replays these exact rational calculations as
an independent finite audit of the analytic formula.
More explicitly, $\Pi$ commutes with $K$, so
$\mathcal O(\Pi(F))\ne0$ implies
$\Pi((F-KF)/2)\ne0$.  The kernel theorem therefore gives
$(F-KF)/2\notin\mathcal J$, which is the asserted imaginary statement.
\end{proof}

The complete $\mathcal O(\Pi(F))$ has $3444$ nonzero word coordinates;
it is provided as \src{data/cayley_s6_normal_form.json}.
The displayed nonzero coefficient alone is sufficient for separation.
This does not contradict the very small certified numerical residual
of $F$.  It proves nonmembership in a specified formal ideal whose
definition omits double shuffle and other available period identities.
In particular, the new certificate does not assert nonmembership in
the separate $5131$-row presentation, does not enlarge that row
presentation by implication, and does not establish or disprove the
numerical conjecture.

\subsection{Verification and the next reduction step}

The independent formulas exposed in the code are useful for auditing
the construction.  The implementation of $e$ uses the equivalent
permutation formula
\[
 e(v_1\cdots v_n)=\sum_{\sigma\in S_n}
 \frac{(-1)^{d(\sigma^{-1})}}
 {n\binom{n-1}{d(\sigma^{-1})}}
 v_{\sigma(1)}\cdots v_{\sigma(n)},
\]
where $d$ counts descents.  It follows directly from
\eqref{cay:eq:eulerian}: a permutation can arise from a cut into $k$
blocks exactly when the $k-1$ cut positions include every descent of
the inverse permutation.  Summing
$(-1)^{k-1}\binom{n-1-d}{k-1-d}/k$ gives the displayed coefficient.
The verifier compares this implementation with the consecutive-cut
definition on small distinct-letter words.  It also checks the
projection and product identities on all words and all shuffle products
of total weight at most four, then reconstructs the weight-seven target
and its separating values with rational arithmetic.

A combined proof search can now replace each additional exact
relation $R$ by $\Pi(R)$ and the target by $\Pi(F)$.  For any specified
family $\mathscr R\subset\mathcal A$, the elementary equivalence
\[
 F\in\mathcal J+\operatorname{span}_{\mathbb Q}\mathscr R
 \quad\Longleftrightarrow\quad
 \Pi(F)\in\operatorname{span}_{\mathbb Q}\{\Pi(R):R\in\mathscr R\}
\]
follows by subtracting $f-\Pi(f)\in\mathcal J$.
This supplies an exact, reconstructible way to impose the full
Cayley ideal before searching the additional double-shuffle rows.
It is a completed reduction tool for that next search, rather than a
claim that the remaining search has already succeeded.
